% Paste these revised sections into the main .tex file supplied by Claude.
% Required preamble packages: \usepackage{booktabs,graphicx,amsmath}

\section{Results}
\label{sec:results}

Table~\ref{tab:final-metrics} reports performance on the held-out test split
($n=624$) at a pre-specified probability threshold of 0.5. Confidence intervals
are percentile-bootstrap 95\% intervals from 2,000 resamples of the test set.
The vision model uses DenseNet-121 with CBAM. The fusion model uses the same
image branch and additionally receives synthetic tabular features; consequently,
the fusion result demonstrates the implementation's ability to use correlated
tabular signal and must not be interpreted as clinical validation.

\begin{table}[htbp]
\centering
\caption{Held-out test performance. Values in brackets are bootstrap 95\%
confidence intervals.}
\label{tab:final-metrics}
\small
\begin{tabular}{lcc}
\toprule
Metric & Vision + CBAM & Fusion + CBAM \\
\midrule
Accuracy  & 0.8638 [0.8365, 0.8895] & 0.8702 [0.8429, 0.8959] \\
Precision & 0.8224 [0.7876, 0.8566] & 0.8280 [0.7927, 0.8627] \\
Recall    & 0.9974 [0.9920, 1.0000] & 1.0000 [1.0000, 1.0000] \\
F1-score  & 0.9015 [0.8800, 0.9217] & 0.9059 [0.8843, 0.9263] \\
AUROC     & 0.9604 [0.9412, 0.9757] & 0.9899 [0.9810, 0.9961] \\
AUPR      & 0.9628 [0.9377, 0.9825] & 0.9921 [0.9839, 0.9978] \\
\bottomrule
\end{tabular}
\end{table}

The fusion configuration has higher point estimates for every reported metric.
Its AUROC increases by 0.0295 and its AUPR by 0.0293 relative to the vision
configuration. The AUROC and AUPR intervals do not overlap; this descriptive
observation is not a formal paired hypothesis test. Threshold-based intervals
for accuracy and F1-score overlap, so the results do not support a strong claim
of improvement for those metrics at the fixed threshold.

\section{Corrections to Experimental Setup}

\subsection{Evaluation Metrics and Uncertainty}
Models are evaluated using accuracy, precision, recall, F1-score, AUROC, and
AUPR. Accuracy, precision, recall, and F1-score are calculated at a fixed
threshold of 0.5; this threshold is not tuned on the held-out test split. AUROC
and AUPR summarize ranking performance across thresholds. For every metric, a
percentile-bootstrap 95\% confidence interval is calculated using 2,000
resamples of the held-out test set with replacement.

\subsection{Data Splitting}
The original 624-image test split is preserved. The original training and
validation pools are merged and then split into training and validation sets
using the best available patient identifiers. Pneumonia filenames encode a
patient identifier; normal-class images do not provide an equivalent identifier
and are therefore treated as unique patients. This reduces, but cannot prove
the complete absence of, patient-level leakage. The resulting experimental
splits contain 4,434 training, 798 validation, and 624 test images.

\subsection{Statistical Replication}
Selected architecture and attention-consistency comparisons are repeated over
three random seeds (42, 123, and 2024) and summarized as mean $\pm$ standard
deviation. Any associated $t$-tests are exploratory because $n=3$ offers low
statistical power; their p-values should not be treated as definitive evidence.

\subsection{Extended Replication: CBAM and Attention-Consistency at Five Seeds}
\label{sec:five-seed}

The three-seed CBAM and attention-consistency comparisons were extended to
five seeds (42, 123, 2024, 7, 2025) to assess whether the exploratory
three-seed estimates were reliable. The two comparisons diverged sharply
under this extension.

For CBAM alone (vision-only, versus no attention), the AUROC difference
shrank from $-0.0122 \pm 0.0139$ ($p=0.31$, three seeds) to
$-0.0037 \pm 0.0179$ ($p=0.67$, five seeds), and the localization
difference shrank from $+0.060 \pm 0.076$ ($p=0.31$) to
$+0.029 \pm 0.107$ ($p=0.57$). One added seed (seed 7) showed a
localization effect in the opposite direction from all other seeds. With
five seeds, CBAM shows no statistically detectable effect on either metric.

For attention-consistency training, extending to five seeds had the
opposite effect: the localization improvement strengthened from
$+0.134 \pm 0.055$ ($p=0.051$) to $+0.083 \pm 0.061$ ($p=0.038$), and the
AUROC cost strengthened from $-0.026 \pm 0.021$ ($p=0.164$) to
$-0.064 \pm 0.032$ ($p=0.012$). Both effects are significant at $\alpha=0.05$
with five seeds. Table~\ref{tab:five-seed} reports the full per-seed values.

\begin{table}[htbp]
\centering
\caption{Five-seed AUROC and lung-localization comparisons. CBAM columns
compare vision-only with and without CBAM. Attention-consistency (AC)
columns compare a CBAM baseline with attention-consistency training added
($\lambda=0.1$).}
\label{tab:five-seed}
\small
\begin{tabular}{lcccc}
\toprule
Seed & CBAM AUROC $\Delta$ & CBAM Loc. $\Delta$ & AC AUROC $\Delta$ & AC Loc. $\Delta$ \\
\midrule
42   & $+0.0016$ & $+0.0955$ & $-0.0375$ & $+0.1143$ \\
123  & $-0.0250$ & $+0.1499$ & $-0.0384$ & $+0.0363$ \\
2024 & $-0.0132$ & $+0.0079$ & $-0.1058$ & $+0.1717$ \\
7    & $+0.0228$ & $-0.1320$ & $-0.0463$ & $+0.0210$ \\
2025 & $-0.0045$ & $+0.0240$ & $-0.0915$ & $+0.0730$ \\
\midrule
Mean $\pm$ SD & $-0.0037 \pm 0.0179$ & $+0.0291 \pm 0.1066$ & $-0.0639 \pm 0.0323$ & $+0.0833 \pm 0.0612$ \\
Paired $t$-test $p$ & $0.67$ & $0.57$ & $0.012$ & $0.038$ \\
\bottomrule
\end{tabular}
\end{table}

While rerunning seeds 123 and 2024 for this analysis, a reproducibility
issue was identified and fixed: \texttt{torch.manual\_seed()} alone did not
make training deterministic in this codebase, because training-time
augmentation draws from Python's \texttt{random} module inside DataLoader
worker processes, which was never separately seeded, and cuDNN's default
convolution algorithms are non-deterministic on GPU. A rerun of seed 2024
under the unfixed code gave a localization value of 0.642 versus an
original run's 0.537 -- a gap larger than the effect size under study. A
seeding utility fixing both issues was added and verified to produce
bit-for-bit identical results across repeated same-seed runs.

\subsection{Calibration}
\label{sec:calibration}

Ranking metrics do not establish whether predicted probabilities are
trustworthy as confidence estimates. Expected Calibration Error (ECE,
equal-width binning, 10 bins) and Brier score were computed on the test
set for both the vision and fusion models. Both models are meaningfully
overconfident: vision ECE $=0.136$ (95\% CI $[0.113, 0.162]$), fusion ECE
$=0.135$ (95\% CI $[0.110, 0.160]$). The miscalibration is concentrated in
the $[0.9,1.0]$ confidence bin, which contains 71\% of test images (442/624)
at a mean stated confidence of 0.995 against an observed accuracy of 0.873.

Three corrective approaches were evaluated, in increasing order of how much
they alter the training or inference pipeline.

\textbf{Temperature scaling}~\cite{guo2017calibration} fits a single scalar
$T$ on validation logits by minimizing cross-entropy, then divides test
logits by $T$ before the sigmoid. This left ECE essentially unchanged
(vision: $0.136 \rightarrow 0.137$; fusion: $0.135 \rightarrow 0.136$) despite
fusion requiring a substantially larger fitted $T$ ($1.21$ versus $1.05$),
consistent with the miscalibration being concentrated in one confidence
region rather than uniformly distributed.

\textbf{Isotonic regression} fits a non-decreasing calibration map on
validation data, capable of correcting different probability ranges
independently. This degraded every metric for the vision model (ECE
$0.136 \rightarrow 0.178$, AUROC $0.960 \rightarrow 0.932$), consistent with
overfitting given only 798 validation examples concentrated away from the
problem region. For fusion, ECE improved ($0.135 \rightarrow 0.113$) but
AUROC fell by $0.041$, too large a cost to adopt.

\textbf{Label smoothing}~\cite{muller2019labelsmoothing} ($\epsilon=0.1$)
was applied during training rather than as a post-hoc correction. This
reduced ECE on both models (vision: $0.136 \rightarrow 0.103$; fusion:
$0.135 \rightarrow 0.112$) with a small apparent AUROC change in each case.
A paired bootstrap (2,000 resamples, shared indices across the
baseline/smoothed comparison) found neither AUROC difference statistically
significant: vision $-0.0143$ (95\% CI $[-0.0309, 0.0023]$, $p=0.087$);
fusion $+0.0028$ (95\% CI $[-0.0032, 0.0104]$, $p=0.42$). Label smoothing is
therefore the only one of the three approaches that improved calibration on
both models without a statistically confirmed cost to ranking performance.

\subsection{Attention Mechanism Comparison: SE}
\label{sec:se}

A Squeeze-and-Excitation (SE) block~\cite{hu2018se} was added to the vision model as an alternative to CBAM, following its use in Potharaju et al.~\cite{potharaju2025enhanced} (discussed below). SE performs channel attention only, via global average pooling and two fully connected layers, without CBAM's added spatial-attention stage. Evaluated across the same five seeds as the CBAM comparison in Section~\ref{sec:five-seed}:

\begin{table}[htbp]
\centering
\caption{Test AUROC and lung-energy localization across attention configurations, five seeds.}
\label{tab:se}
\small
\begin{tabular}{lcc}
\toprule
Configuration & AUROC (mean $\pm$ SD) & Localization (mean $\pm$ SD) \\
\midrule
No attention & 0.9595 $\pm$ 0.0184 & 0.439 $\pm$ 0.026 \\
CBAM & 0.9559 $\pm$ 0.0079 & 0.468 $\pm$ 0.086 \\
SE & 0.9632 $\pm$ 0.0089 & 0.479 $\pm$ 0.049 \\
\bottomrule
\end{tabular}
\end{table}

SE has the best mean of the three configurations on both metrics, but neither difference from CBAM is significant (AUROC: $p=0.099$; localization: $p=0.84$), nor is either difference from no attention (AUROC: $p=0.72$; localization: $p=0.16$). This is read as a mild, real-looking, but statistically unconfirmed effect, consistent with the pattern established for CBAM alone above: single-digit-seed comparisons on this dataset do not reliably separate small effects from noise.

\subsection{External Comparison: Split-Sensitivity Analysis}
\label{sec:potharaju}

Potharaju et al.~\cite{potharaju2025enhanced} report 98.6\% accuracy for a CBAM-enhanced CNN on what their cited data source confirms is the same underlying dataset used in this work. Their own results table is internally inconsistent on this exact result: reported precision (98.4\%) and recall (98.3\%) are incompatible with their reported F1 (94.5\%), which should also be approximately 98.3\%. Their abstract and conclusion also disagree with each other on a separate result (SE+CNN accuracy reported as 96.25\% and 96.17\% respectively). Their paper does not describe patient-level splitting, despite this dataset's pneumonia-class filenames encoding a shared patient identifier across multiple images, and their baseline CNN's architecture (layer count, filter sizes, dense layer width) is not specified anywhere in the paper.

Reproducing their reported split sizes (5216 train / 160 validation / 480 test, stratified by class) found that 304 of 480 test images (63\%) share a patient identifier with an image in the training set. Training this work's existing DenseNet-121+CBAM model, unchanged, on this split across three seeds gave the results in Table~\ref{tab:potharaju}.

\begin{table}[htbp]
\centering
\caption{This work's unmodified model under Potharaju et al.'s split protocol, compared with their reported result and with this work's own patient-level split.}
\label{tab:potharaju}
\small
\begin{tabular}{lccccc}
\toprule
Setting & Accuracy & Precision & Recall & F1 & AUROC \\
\midrule
Potharaju et al. (as reported) & 98.6\% & 98.4\% & 98.3\% & 94.5\%$^*$ & not reported \\
Their split, seed 42 & 96.88\% & 98.83\% & 96.86\% & 97.84\% & 0.9969 \\
Their split, seed 123 & 98.33\% & 98.03\% & 99.71\% & 98.87\% & 0.9984 \\
Their split, seed 2024 & 96.46\% & 98.54\% & 96.57\% & 97.55\% & 0.9967 \\
Their split, 3-seed mean & 97.22\% & 98.47\% & 97.71\% & 98.09\% & 0.9973 \\
This work's own split & 86.38\% & 82.24\% & 99.74\% & 90.15\% & 0.9604 \\
\bottomrule
\end{tabular}
\end{table}

\noindent $^*$Internally inconsistent given the reported precision and recall; see above.

Under their split protocol, this work's unmodified model reaches a mean accuracy of 97.22\% (best seed 98.33\%), in the same range as their reported 98.6\% but not exceeding it, and with no architecture changes from the model evaluated at 86.38\% on this work's own patient-level split. This does not establish what their actual pipeline did, since no code was released; it shows that a split matching their description can produce performance from an unmodified model in the range of their reported result. The large difference between the two rows of this work's own results (97.22\% vs. 86.38\%, same model) is therefore consistent with split sensitivity being a major contributor to the observed performance gap.

As a second, independent test, SE (Section~\ref{sec:se}) -- Potharaju et al.'s one precisely specified architectural contribution -- was evaluated on this work's own honest split rather than on theirs, isolating the architecture question from the split question. It gave a small, non-significant improvement over CBAM (Table~\ref{tab:se}), far short of the gap to their reported number, reinforcing that the split is the larger factor.

% Replace the existing Proposed row in Table 1 with this row. Do not use the
% earlier 0.9608 value: the evaluated checkpoint has AUROC 0.9604.
% \textbf{Proposed} & \textbf{DenseNet121 + CBAM} & \textbf{Binary CXR}
% & \textbf{86.38\%} & \textbf{82.24\%} & \textbf{99.74\%}
% & \textbf{90.15\%} & \textbf{0.9628} & \textbf{0.9604} \\
